\documentclass[10pt,a4paper]{amsart}
\usepackage{amsmath,amssymb,amsthm,graphicx}
% Retained minimal preamble for the Stacks Project chapter varieties.tex.
% Verbatim excerpt of the chapter preamble: Xy-pic setup (native preamble lines 14-18), the native theorem
% declarations with the mathematical macros (native preamble lines 154-217), and the project own no-op
% definitions of the reference/slogan/history display environments (native wrapper.tex lines 1-3) that the
% counted statement uses for its slogan. Copied mechanically; nothing added, removed or redefined.
% For commutative diagrams we use Xy-pic
\usepackage[all]{xy}

% We use 2cell for 2-commutative diagrams.
\xyoption{2cell}
% Theorem environments.
%
\theoremstyle{plain}
\newtheorem{theorem}[subsection]{Theorem}
\newtheorem{proposition}[subsection]{Proposition}
\newtheorem{lemma}[subsection]{Lemma}

\theoremstyle{definition}
\newtheorem{definition}[subsection]{Definition}
\newtheorem{example}[subsection]{Example}
\newtheorem{exercise}[subsection]{Exercise}
\newtheorem{situation}[subsection]{Situation}

\theoremstyle{remark}
\newtheorem{remark}[subsection]{Remark}
\newtheorem{remarks}[subsection]{Remarks}

\numberwithin{equation}{subsection}

% Macros
%
\def\lim{\mathop{\mathrm{lim}}\nolimits}
\def\colim{\mathop{\mathrm{colim}}\nolimits}
\def\Spec{\mathop{\mathrm{Spec}}}
\def\Hom{\mathop{\mathrm{Hom}}\nolimits}
\def\Ext{\mathop{\mathrm{Ext}}\nolimits}
\def\SheafHom{\mathop{\mathcal{H}\!\mathit{om}}\nolimits}
\def\SheafExt{\mathop{\mathcal{E}\!\mathit{xt}}\nolimits}
\def\Sch{\mathit{Sch}}
\def\Mor{\mathop{\mathrm{Mor}}\nolimits}
\def\Ob{\mathop{\mathrm{Ob}}\nolimits}
\def\Sh{\mathop{\mathit{Sh}}\nolimits}
\def\NL{\mathop{N\!L}\nolimits}
\def\CH{\mathop{\mathrm{CH}}\nolimits}
\def\proetale{{pro\text{-}\acute{e}tale}}
\def\etale{{\acute{e}tale}}
\def\QCoh{\mathit{QCoh}}
\def\Ker{\mathop{\mathrm{Ker}}}
\def\Im{\mathop{\mathrm{Im}}}
\def\Coker{\mathop{\mathrm{Coker}}}
\def\Coim{\mathop{\mathrm{Coim}}}

% Boxtimes
%
\DeclareMathSymbol{\boxtimes}{\mathbin}{AMSa}{"02}

%
% Macros for moduli stacks/spaces
%
\def\QCohstack{\mathcal{QC}\!\mathit{oh}}
\def\Cohstack{\mathcal{C}\!\mathit{oh}}
\def\Spacesstack{\mathcal{S}\!\mathit{paces}}
\def\Quotfunctor{\mathrm{Quot}}
\def\Hilbfunctor{\mathrm{Hilb}}
\def\Curvesstack{\mathcal{C}\!\mathit{urves}}
\def\Polarizedstack{\mathcal{P}\!\mathit{olarized}}
\def\Complexesstack{\mathcal{C}\!\mathit{omplexes}}
% \Pic is the operator that assigns to X its picard group, usage \Pic(X)
% \Picardstack_{X/B} denotes the Picard stack of X over B
% \Picardfunctor_{X/B} denotes the Picard functor of X over B
\def\Pic{\mathop{\mathrm{Pic}}\nolimits}
\def\Picardstack{\mathcal{P}\!\mathit{ic}}
\def\Picardfunctor{\mathrm{Pic}}
\def\Deformationcategory{\mathcal{D}\!\mathit{ef}}
\newenvironment{reference}{}{}
\newenvironment{slogan}{}{}
\newenvironment{history}{}{}
\usepackage{cleveref}
\crefname{originallemma}{original lemma}{original lemmas}
\Crefname{originallemma}{Original lemma}{Original lemmas}
\makeatletter
\newlabel{coherent-lemma-flat-base-change-cohomology}{{02KH}{source}{}{sourcepointer-1}{}}
\newlabel{coherent-lemma-flat-base-change-cohomology@cref}{{[originallemma][0][]02KH}{[1][0][]source}{}{}{}}
\makeatother
\title{Global generation of a quasi-coherent module descends along a faithfully flat ring map}
\author{The Stacks Project Authors}
\date{stacks-project:varieties.tex:local-snapshot-20261011; GNU FDL 1.2 or later}
\begin{document}
\maketitle
\noindent\textit{Editorial note.} All delivered mathematical text is the author's own, copied from the pinned Stacks Project chapter file varieties.tex; no statement or proof step was rewritten, reordered or removed.\par\medskip
\noindent\textit{Editorial note.} The proof is complete as written and ends with the faithfully flat criterion; no step is referred to the reader.\par\medskip
\noindent\textit{Editorial note.} The counted spans carry 1 navigation stub and no bibliography citation.\par\medskip
\noindent\textit{Editorial note.} The preamble is the chapter's retained minimal preamble (sources/stacks-varieties/retained-preamble.tex lines 1-77), a verbatim excerpt of the native preamble.\par\medskip
\noindent\textit{Editorial note.} The label lemma-globally-generated-base-change is defined only in varieties.tex, so no cross-chapter restatement had to be weighed for this card.\par\medskip
\noindent\textit{Allowed standard prerequisites (editorial).}\begin{itemize}
\item Flat base change for the cohomology of a quasi-coherent module. Cohomology of Schemes, Lemma coherent-lemma-flat-base-change-cohomology, invoked at proof lines 4213-4215. Supplies the identification of the global sections of the base changed module with the tensor product of the global sections.
\item Surjectivity of a module map is equivalent to surjectivity after a faithfully flat base change. Standard algebra, used at proof lines 4226-4229. Supplies the equivalence for the evaluation map on an affine chart.
\end{itemize}\medskip
\ifdefined\hypertarget\hypertarget{sourcepointer-1}{}\fi
\noindent\textit{Source reference.} Original lemma 02KH: sources/stacks-chow/tags.txt, source lines 3372--3372. Pointer outside the counted span, invoked once for context; the load-bearing argument of the counted statement is the author's proof reproduced in full in the proof span.\par\medskip
\begin{lemma}
\label{lemma-globally-generated-base-change}
Let $X \to \Spec(A)$ be a morphism of schemes. Let $A \subset A'$
be a faithfully flat ring map. Let $\mathcal{F}$ be a quasi-coherent
$\mathcal{O}_X$-module. Then $\mathcal{F}$ is globally generated
if and only if the base change $\mathcal{F}_{A'}$ is globally generated.
\end{lemma}

\begin{proof}
More precisely, set $X_{A'} = X \times_{\Spec(A)} \Spec(A')$.
Let $\mathcal{F}_{A'} = p^*\mathcal{F}$ where $p : X_{A'} \to X$
is the projection. By
Cohomology of Schemes, Lemma \ref{coherent-lemma-flat-base-change-cohomology}
we have
$H^0(X_{k'}, \mathcal{F}_{A'}) = H^0(X, \mathcal{F}) \otimes_A A'$.
Thus if $s_i$, $i \in I$ are generators for $H^0(X, \mathcal{F})$
as an $A$-module, then their images in $H^0(X_{A'}, \mathcal{F}_{A'})$
are generators for $H^0(X_{A'}, \mathcal{F}_{A'})$ as an $A'$-module.
Thus we have to show that the map
$\alpha : \bigoplus_{i \in I} \mathcal{O}_X \to \mathcal{F}$,
$(f_i) \mapsto \sum f_i s_i$
is surjective if and only if $p^*\alpha$ is surjective.
This we may check over an affine open $U = \Spec(B)$ of $X$.
Then $\mathcal{F}|_U$ corresponds to a $B$-module $M$
and $s_i|_U$ to elements $x_i \in M$. Thus we have to show
that $\bigoplus_{i \in I} B \to M$ is surjective if and only
if the base change $\bigoplus_{i \in I} B \otimes_A A' \to M \otimes_A A'$
is surjective. This is true because $A \to A'$ is faithfully flat.
\end{proof}

\end{document}
